\section{Results and Interpretation}

\subsection{Baseline Results: Pressure and Post-Earnings Returns}

We begin by examining whether the microstructure-based pressure measure predicts post-earnings returns. Table~\ref{tab:robust_baseline} reports cross-sectional regressions of 20-day forward returns on the pressure score across alternative constructions.

Across specifications, we find a consistent negative relationship between pressure and subsequent returns. In particular, methods A, B, and D yield negative and statistically significant coefficients under firm, date, and two-way clustered standard errors. Method D produces the strongest result, with a standardized coefficient of $\beta^{\text{std}} = -0.00595$ and a two-way clustered $t$-statistic of $-3.162$.

Economically, this implies that a one-standard-deviation increase in pressure is associated with approximately 60 basis points lower returns over the subsequent 20 trading days. While the explanatory power is modest, the magnitude and statistical significance are consistent with typical results in cross-sectional return prediction.

The negative relationship suggests that higher values of the pressure score—corresponding to more information-dominated trading—are associated with attenuated post-earnings-announcement drift. This is consistent with a pre-announcement information incorporation mechanism, whereby informed trading prior to earnings reduces the amount of information remaining to be incorporated after the announcement.

\subsection{Interaction with Earnings Surprise}

We next examine whether the predictive content of the pressure score is distinct from the information contained in earnings surprises. Table~\ref{tab:robust_sue} reports regressions including standardized unexpected earnings (SUE) and interaction terms.

Consistent with the PEAD literature, the coefficient on SUE is strongly positive across all specifications, confirming that earnings surprises drive post-announcement return continuation. Importantly, the pressure score retains its negative sign and remains statistically significant for the stronger constructions after controlling for SUE.

To further isolate the role of microstructure-based factors, we estimate specifications including both PCA-derived factors and SUE simultaneously. The results confirm that the information factor remains negative and statistically significant ($\beta = -0.0045$, $t = -3.12$) even after controlling for SUE, while the SUE coefficient is strongly positive ($\beta = 0.0064$, $t = 4.61$).

The interaction between pressure and SUE is negative and significant, indicating that the attenuation effect of pressure is stronger when earnings surprises are larger. This suggests that pre-announcement trading pressure plays a more important role precisely in events with greater fundamental information.

Taken together, these findings imply that the pressure score does not simply proxy for earnings news, but instead captures the timing and mechanism of information incorporation. While SUE measures the magnitude of fundamental news, microstructure-based measures reflect how and when that information is incorporated into prices.
\begin{table}[!t]
\centering
\caption{Regression with SUE Control}
\label{tab:sue_control}
\begin{tabular}{lccc}
\hline
 & Coefficient & t-stat & Significance \\
\hline
InfoFactor & -0.0045 & -3.12 & *** \\
LiquidityFactor & -0.0024 & -2.16 & ** \\
SUE & 0.0064 & 4.61 & *** \\
Constant & 0.0139 & 9.66 & *** \\
\hline
Observations & \multicolumn{3}{c}{4127} \\
$R^2$ & \multicolumn{3}{c}{0.011} \\
\hline
\end{tabular}
\end{table}

\subsection{Latent Structure: PCA Evidence}

To further investigate the structure of the microstructure feature space, we apply principal component analysis (PCA) to the standardized event-level features. The first three principal components explain approximately 54\% of the total variance, indicating that the feature space is characterized by multiple economically meaningful dimensions.

Inspection of the component loadings reveals distinct patterns corresponding to different trading mechanisms. In particular, the fifth principal component (PC5) loads strongly on signed order flow measures, including mean OFI and OFI autocorrelation. These features capture both the direction and persistence of trading and are consistent with theoretical signatures of informed trading. We therefore interpret PC5 as an information-related component.

Importantly, the fact that the information-related component emerges as a higher-order principal component rather than the first component highlights that directional and persistent order flow is not the dominant source of variation in the data, but instead represents a more subtle and economically meaningful dimension of trading behavior.

By contrast, the second principal component (PC2) loads heavily on spread-based measures, including spread volatility and tail behavior, indicating sensitivity to liquidity conditions and trading frictions. The third component (PC3) loads on absolute order flow and trading activity measures, capturing the intensity and magnitude of trading. Together, PC2 and PC3 reflect dimensions of liquidity stress and non-directional trading pressure.

Based on these patterns, we construct two PCA-based factors:
\begin{equation}
\text{InfoFactor}_i = \text{PC5}_i, \qquad
\text{LiquidityFactor}_i = \text{PC2}_i + \text{PC3}_i.
\end{equation}

\subsection{PCA-Based Predictive Regressions}

We next evaluate whether the PCA-derived factors have predictive content for post-earnings returns. Table~\ref{tab:pca_regression} reports regressions of 20-day forward returns on the information and liquidity factors.

The information factor exhibits a negative and statistically significant coefficient ($\beta = -0.0045$, $t = -3.09$), indicating that higher values of the information-related component are associated with lower subsequent returns. This result closely mirrors the baseline pressure findings and reinforces the interpretation that persistent directional order flow reflects pre-announcement information incorporation.

The liquidity factor also enters with a negative coefficient ($\beta = -0.0026$, $t = -2.34$), although its magnitude is smaller. This suggests that liquidity-related trading pressure may also contribute to return predictability, but plays a secondary role relative to information-driven trading.

Although the overall explanatory power remains modest ($R^2 = 0.006$), the joint significance of the factors indicates that latent microstructure components contain systematic information about future returns.

To further illustrate the economic relevance of the information factor, we sort events into quintiles based on $\text{InfoFactor}_i$ and compute average post-event returns. The resulting portfolio returns display a clear decline at higher quintiles, with the highest information-factor quintile yielding an average return of approximately 0.45\%, compared to 1.54\% for the lowest quintile.

The spread between the highest and lowest quintiles (Q5--Q1) is approximately $-1.10\%$ over the 20-day post-event window. While the intermediate quintiles are not perfectly monotonic, the concentration of the effect in the upper tail suggests that strong and persistent directional order flow is required to generate economically meaningful predictability.
\begin{table}[!t]
\centering
\caption{PCA-Based Predictive Regression}
\label{tab:pca_regression}
\begin{tabular}{lccc}
\hline
 & Coefficient & t-stat & Significance \\
\hline
InfoFactor & -0.0045 & -3.09 & *** \\
LiquidityFactor & -0.0026 & -2.34 & ** \\
Constant & 0.0140 & 9.74 & *** \\
\hline
Observations & \multicolumn{3}{c}{4144} \\
$R^2$ & \multicolumn{3}{c}{0.006} \\
\hline
\end{tabular}
\end{table}

Importantly, the PCA-based factors and the economically constructed pressure score capture closely related dimensions of trading behavior. The alignment between the pressure score results and the PCA-based information factor provides additional validation that the reduced-form construction successfully isolates information-driven trading pressure.

\subsection{Cross-Sectional Heterogeneity: Size Effects}

We next examine whether the predictive power of the PCA-based factors varies across firm size. Regressions estimated separately for small-cap and large-cap subsamples reveal that the information factor has a stronger effect among smaller firms.

In the small-cap subsample, the coefficient on the information factor is $-0.0055$ ($t = -2.30$), compared to $-0.0038$ ($t = -2.24$) for large-cap firms. In addition, the explanatory power is higher for small-cap stocks ($R^2 = 0.015$ versus $0.007$).

The liquidity factor also exhibits a stronger (though less precisely estimated) effect in small-cap stocks, while it is weaker and statistically insignificant in large-cap stocks. This pattern is consistent with the idea that liquidity frictions are more pronounced in less liquid securities.

These findings align with market microstructure theory, which predicts that smaller firms exhibit higher information asymmetry and slower price discovery. As a result, informed trading is more likely to leave detectable signatures in order flow and to have a persistent impact on prices.
\begin{table}[!t]
\centering
\caption{Size-Based Heterogeneity}
\label{tab:size_split}
\begin{tabular}{lcc}
\hline
 & Small Cap & Large Cap \\
\hline
InfoFactor & -0.0055** & -0.0038** \\
LiquidityFactor & -0.0037* & -0.0016 \\
SUE & 0.0081*** & 0.0044*** \\
$R^2$ & 0.015 & 0.007 \\
Observations & 2062 & 2065 \\
\hline
\end{tabular}
\end{table}

\subsection{Economic Interpretation}

Taken together, the results suggest a coherent economic mechanism linking pre-announcement trading to post-announcement returns. The baseline pressure score, PCA-derived information factor, and interaction with earnings surprises all point to the same conclusion: persistent directional order flow prior to earnings announcements is associated with attenuated subsequent drift.

This pattern is consistent with models of informed trading in which private information is gradually incorporated into prices through order flow. When information is partially revealed before the announcement, the remaining price adjustment after the announcement is reduced, leading to weaker drift.

At the same time, the presence of liquidity-related components highlights that not all trading pressure is informational. While both information-driven and liquidity-driven trading contribute to return predictability, the evidence indicates that persistent directional order flow plays the dominant role.

Overall, these findings suggest that microstructure-based measures capture economically meaningful variation in the timing and mechanism of information incorporation, complementing traditional earnings-based signals in explaining post-earnings return dynamics.

